\section*{NeurIPS Paper Checklist}

\begin{enumerate}
\item \textbf{Claims.} \answerYes{} The abstract, Introduction, claim card, Results, and Limitations explicitly scope the claims to matched-budget diagnostic evaluation, quality-conditioned evaluation claims, and explicit non-claims.

\item \textbf{Limitations.} \answerYes{} Section 8 discusses the narrow task setting, matched-budget boundary, code-aware episode construction, same-test diagnostic selection, retrospective prior-quality audit, prompt-only negative result, and external/backend sensitivity.

\item \textbf{Theory assumptions and proofs.} \answerNA{} The paper does not present theoretical results or formal proofs.

\item \textbf{Experimental result reproducibility.} \answerYes{} The paper describes the benchmark setting, seeds, candidate and execution budgets, selector, information access, prompt-only controls, and artifact scripts. The artifact regenerates paper-facing tables from cached outputs and verifies provenance.

\item \textbf{Open access to data and code.} \answerYes{} An anonymized artifact is submitted with code, cached outputs, documentation, claim-to-evidence mapping, and reviewer quickstart scripts. Public release will follow upstream redistribution constraints for benchmark-derived resources.

\item \textbf{Experimental setting/details.} \answerYes{} Sections 3--5 specify the task setting, candidate budget, execution budget, seeds, backend, retrieval regimes, selector status, and information-access boundaries. Additional details are included in the artifact.

\item \textbf{Experiment statistical significance.} \answerYes{} The paper reports mean and standard deviation over seeds, paired query-seed directionality, and task-clustered directionality. It avoids overclaiming statistical independence of query-seed outcomes.

\item \textbf{Experiments compute resources.} \answerYes{} The artifact documents the quickstart environment and distinguishes CPU-only cached-output regeneration from optional live LLM reruns requiring matching model endpoints or local models. The main paper reports that prompt tokens are not matched and that live reruns are separate from cached-output verification.

\item \textbf{Code of ethics.} \answerYes{} The authors have considered the NeurIPS Code of Ethics. The submission makes no deployment or safety claim and warns that generated-code execution during live reruns should be sandboxed.

\item \textbf{Broader impacts.} \answerYes{} The work is an evaluation audit rather than a deployed system. Potential risks are mainly around misinterpreting diagnostic evidence as deployable evidence or executing generated code unsafely; these are discussed in Sections 7--8.

\item \textbf{Safeguards.} \answerNA{} The paper does not release a pretrained model or high-risk generative model. The artifact distinguishes safe cached-output verification from optional live generated-code execution, which requires sandboxing.

\item \textbf{Licenses.} \answerYes{} The paper cites the benchmark and model sources used in the evaluation. The artifact includes source/data and release-policy documentation and avoids public redistribution of benchmark-derived resources unless upstream permission is verified.

\item \textbf{Assets.} \answerYes{} The submitted artifact documents its code, cached outputs, provenance, information-access card, known limitations, and reproducibility status. It is anonymized for review.

\item \textbf{Crowdsourcing and research with human subjects.} \answerNA{} The paper does not involve crowdsourcing or human-subject data collection.

\item \textbf{IRB approvals.} \answerNA{} The paper does not involve human subjects or IRB-regulated data collection.

\item \textbf{Declaration of LLM usage.} \answerNA{} LLM tools were used only for writing, editing, organization, and claim-risk review, not as an important, original, or non-standard component of the core methodology. All claims, citations, numerical results, tables, artifact descriptions, and limitations were checked and approved by the human authors, who take full responsibility for the submission.
\end{enumerate}
